\subsection{All nonunit odd terminals: choosing the actual eighth-moment direction}\label{proof:R2-C030}
Let the strict odd residual list at the terminal be
\fitdisplay{
\Sigma=(\ell,r_1,t_1,\ldots,r_j,t_j),\qquad
 2\le\ell<r_1<t_1<\cdots<r_j<t_j,\quad j\ge1.
}{R30.1.1}
Keep its genuine equal-pair product $E_t$. For the existing receiving
arm $m\ge3$ and connector $L\ge1$, put $n=\ell+L$ and choose
\fitdisplay{
F_\Sigma=\begin{cases}
 \bar f_\ell\prod_{i=1}^j h_{r_i}f_{t_i},&r_1\ge4,\\
 \bar f_2f_3h_c\prod_{i=2}^jh_{r_i}f_{t_i},&r_1=3,\ \ell=2,\ c=t_1\ge4.
 \end{cases}
}{R30.1.2}
Empty products equal one. If $m\le n$, the receiving arm is
$W_m=A\bar f_{m-1}$ and the connector is $T^L$.
If $m>n$, the receiving arm is $\bar f_m=T^m\one$ and the connector is
$AT^{L-1}$. Thus the first connector edge changes only according to
the actual integers. When $L=1$ the latter operator is $A$ on the
one existing edge. Denote the chosen receiving arm by $R_{m,n}$ and
the connector by $\mathcal P_{m,n,L}$. We prove
\fitdisplay{
\E[\{R_{m,n}\mathcal P_{m,n,L}(E_tF_\Sigma)\}^{8}]
 \le e^{-3d}.
}{R30.1.3}
The same bound holds after a genuine equal-pair product $E_c$ and
any finite ordered negative post block
$P_c^-=\prod_{\nu=1}^kW_{u_\nu}\bar f_{v_\nu}$,
$4\le u_1<v_1<\cdots<u_k<v_k$, $k\ge0$, are restored at the receiver.
The terminal and receiver length intervals may overlap. This is a
construction for the stated direction.

\paragraph{R30.2 Finite tools}
Use $B=-K$, $u=B\one=-g$, $T=J+B$, $A=J-B$, $p=1-u$,
$q=1+u$, $\delta_i=\|B^i\one\|_2^2$, $d=\delta_1$,
$\theta=13/32$, $\kappa=19/32$ and $\tau=d-\delta_2$.
The finite spectral and row bounds give
\fitdisplay{
d^2\le\delta_2\le\theta d,\quad d\le1/3,\quad
 \delta_{i+1}\le\theta\delta_i,\quad\|B\|^2\le\theta,
 \quad |BV(x)|\le\|V\|_2.
}{}
For $G_s(z)=\sum_{i=0}^{s-1}z^i$, $G_0=0$, one has
$\|G_s(B)\|\le1$ for $s=0,1$ or $s\ge3$, while
$\|G_2(B)\|^2\le1+\theta$. Define the finite sums
\fitdisplay{
\beta_1=\beta_2=0,\qquad
 \beta_n=\sum_{i=1}^{\lfloor(n-1)/2\rfloor}(-1)^{i-1}\delta_i,
 \qquad \tau\le\beta_n\le d\quad(n\ge3).
}{}
The spectral measure $\nu$ at $u$ has mass $d$, support
$[0,\theta]$ and moments $\int x^i\,d\nu=\delta_{i+1}$.
Writing $z=i\sqrt x$, a real polynomial $P=P_e+zP_o$ satisfies
$\E P(B)u=-\int P_o\,d\nu$. All four path messages lie in $[0,2]$;
their means satisfy $\mu_a=\E f_a=\E\bar f_a\le1$,
$\mu_2=1-d$, $\mu_a\le1-2d$ for $a\ge3$, and
$W_a+\bar f_a=2\mu_{a-1}$. Only genuine equal pairs are bounded by
one pointwise.

The full-mean exponential lemma of Section~\ref{proof:Q5-C081} is
\fitdisplay{
\E e^{8BV}\le\exp\{8\E(BV)+17\|V\|_2^2\}.
}{R30.2.5}
It applies to every real $V$, including its constant component.
The actual pair-sum estimates of Section~\ref{proof:Q5-C090}, after
full reversal, are
\fitdisplay{
W_h+\bar f_H\le2+\sum_{a=h}^{H-1}B^au-\gamma_h,
 \quad 3\le h<H,\qquad
 \gamma_h=\begin{cases}2d-\delta_2,&h=3,\\16d/5,&h\ge4.
 \end{cases}
}{}
For $4\le r<t$, the actual positive terminal pair satisfies
$h_rf_t\le\exp\{\sum_{a=r}^{t-1}K^ag-16d/5\}$.
A finite disjoint-interval mask with least power $a$ has bounds
\fitdisplay{
\|R\|_2^2\le\theta^a C_*d,\quad
 |R|^2\le\theta^{a-1}C_*d,\qquad C_*=32768/16245.
}{}
These bounds apply to the finite sums in the path expansions.

\paragraph{R30.3 Matching the two actual lengths}
For $3\le h\le H$, put
\fitdisplay{
Q_{h,H}=W_h\bar f_H,\quad
 R_* =\sum_{a=h}^{H-1}B^au=BZ_*,\quad
 Z_*=B^{h-1}G_{H-h}(B)u.
}{}
At $h=H$ these sums vanish. Uniformly,
\fitdisplay{
W_h+\bar f_H\le2-\lambda d+R_*,\quad
 \lambda=51/32,\quad \E R_*\le\theta d,\quad
 \|Z_*\|_2\le\tfrac12\sqrt d.
}{R30.3.2}
For $h<H$ the first assertion follows from the pair-sum estimate,
since $2d-\delta_2\ge51d/32$. At equality use
$2\mu_{h-1}\le2-2d$, including $h=3$. In the finite alternating
mean sum the first possible positive term is $\delta_2$, so
$\E R_*\le\delta_2\le\theta d$. Furthermore
\fitdisplay{
\|Z_*\|_2^2\le\theta^{h-1}(1+\theta)d
 \le7605d/32768<d/4,
}{}
including the $G_2$ exception. Both actual arms are nonnegative.
Their arithmetic-geometric mean bound and $1+x\le e^x$ give
$Q_{h,H}\le e^{R_*-\lambda d}$. Applying (R30.2.5) once yields
\fitdisplay{
\E Q_{h,H}^8\le
 \exp\{(-8\lambda+8\theta+17/4)d\}=e^{-21d/4}.
}{R30.3.4}

\paragraph{R30.4 Both first-edge choices and every finite propagation length}
For $\sigma\in\{+1,-1\}$ let
\fitdisplay{
S_\sigma=J+\sigma B,\quad P_L^\sigma=S_\sigma T^{L-1},\quad
 w_0^\sigma=1,\quad w_j^+=\bar f_j,\quad w_j^-=W_j\ (j\ge1).
}{}
Choose $\sigma=+$ when $m\le n$, and $-$ otherwise.
Suppose $F\le\bar f_\ell e^{R-\Gamma}$ and write
$D_i=\E B^iR$. For the actual range $|R|\le1/2$ verified below,
\eqref{eq:oc:remainder} with $t=1/2$, $c=6/5$ gives
$e^R\le1+R+3R^2/5$.

If $r_1\ge4$, the finite first row is
\fitdisplay{
S_\sigma(\bar f_\ell e^R)
 \le w_{\ell+1}^\sigma+\sigma\chi BR+\delta,
 \quad \chi=\mu_\ell\in[0,1],
}{}
\fitdisplay{
\delta=\mu_\ell D_0+
 2\operatorname{sd}(\bar f_\ell)\|R\|_2+\tfrac{12}{5}\|R\|_2^2.
}{R30.4.2}
Indeed split $S_\sigma(\bar f_\ell R)$ into
$\mu_\ell D_0+\sigma\mu_\ell BR$ and the centered part, which is
at most $2\E| (\bar f_\ell-\mu_\ell)R|$. The two positive kernels
have the same absolute row bound, and the quadratic fee is
$12\E R^2/5$. Thus the same signed $\delta$ works for both first
edges.

If $r_1=3$, the actual first row instead has $\chi=1$:
\fitdisplay{
S_\sigma(\bar f_2e^R)\le w_3^\sigma+\sigma BR+\delta,
}{}
\fitdisplay{
\delta=\E R+\E(uR)+\E((Bu)R)
       +\sqrt{(d+\delta_2)\E R^2}+\tfrac{12}{5}\E R^2.
}{R30.4.4}
Here $\bar f_2=1+u+Bu$; the displayed three means are precisely
$J(\bar f_2R)$, and the remaining row term is bounded by
$\E|(u+Bu)R|\le\sqrt{(d+\delta_2)\E R^2}$.
This proves sign compatibility without identifying the two distinct
first-row coefficients.

For any $D\ge\delta$ and every finite $L\ge1$, with $n=\ell+L$,
\fitdisplay{
e^\Gamma P_L^\sigma F\le
 Y_L^\sigma=w_n^\sigma+\sigma\chi B^LR
       +Dw_{L-1}^\sigma
       +\chi\sum_{i=1}^{L-1}D_iw_{L-1-i}^\sigma.
}{R30.4.5}
For $L=1$ this is the corresponding $S_\sigma$ first row. For $L\ge2$,
apply $S_\sigma T^{L-2}$ to the $T$ first row and use
\eqref{eq:oc:finite-duhamel}. Its final summand is
$\chi D_{L-1}w_0^\sigma$; all indices are nonnegative and positivity gives
$Y_L^\sigma\ge0$.
Replacing signed coefficients by positive parts only when they
multiply a nonnegative message gives the scalar upper bounds
\fitdisplay{
\begin{array}{c|c}
 L& C\ \text{bounding }Dw_{L-1}^\sigma+
 \chi\sum_{i=1}^{L-1}(D_i)_+w_{L-1-i}^\sigma\\\hline
 1&D\\2&2D+(D_1)_+\\3&2D+2(D_1)_++(D_2)_+\\
 L\ge4&2D+2\sum_{i\ge1}(D_i)_+
 \end{array}
}{R30.4.6}
The last geometric bound dominates each finite positive subsum.

The product of the actual receiver and $w_n^\sigma$ is always
$Q_{h,H}$ with $h=\min(m,n)$ and $H=\max(m,n)$.
Consequently
\fitdisplay{
R_{m,n}+Y_L^\sigma\le2-\lambda d+C+BV,
 \qquad V=Z_*+\sigma\chi B^{L-1}R.
}{}
If $\|B^{L-1}R\|_2\le\rho_L\sqrt d$, then
$|D_L|\le\rho_Ld$, and hence
$\E BV\le(\theta+\rho_L)d$ and
$\|V\|_2\le(1/2+\rho_L)\sqrt d$.
Arithmetic-geometric mean followed by (R30.2.5) gives the complete
same-kernel payment
\fitdisplay{
\begin{aligned}
 \log\E[\{R_{m,n}P_L^\sigma F\}^8]
 \le{}&-8\Gamma-\tfrac{51}{4}d+8C
       +8(\theta+\rho_L)d+17(1/2+\rho_L)^2d.
 \end{aligned}
}{R30.4.9}
If the moment is zero use the equivalent exponential inequality.
The terminal mask and root interval need not be disjoint: this is one
exponential expectation, with every scalar cost displayed.

\paragraph{R30.5 Second-shortest terminal length at least four}\label{part:R30.5}
Section~\ref{proof:Q5-C181} gives, for the first row of (R30.1.2),
\fitdisplay{
\begin{gathered}
 R=\sum_{i=1}^j\sum_{a=r_i}^{t_i-1}K^ag,\quad
 F_\Sigma\le\bar f_\ell e^{R-\Gamma},\\
 \Gamma\ge16jd/5,\quad \|R\|_2<\tfrac14\sqrt d,
 \quad |R|^2<d/3,
\end{gathered}
}{}
\fitdisplay{
\delta\le5d/6,\qquad\sum_{i\ge1}(D_i)_+<2d/5.
}{R30.5.1}
For $j=1,r_1=4$ the stronger three entries are
$\sqrt d/5,72d/125,7d/25$; for $j=1,r_1\ge5$ they are
$3\sqrt d/20,3d/5,d/5$. For $j\ge2$ use the joint row, without
combining incompatible extrema. The finite mask of least original
power four can be written $R=(U_e+zU_o)u$ with
\fitdisplay{
-\frac{x^2}{1-x^2}\le U_e\le\frac{x^3}{1-x^2},\qquad
 -\frac{x^3}{1-x^2}\le U_o\le\frac{x^2}{1-x^2}.
}{}
Its four residue classes give the positive-tail bound
\fitdisplay{
\frac{x^2(1+x)(1+x^2)}{(1-x^2)^2}
 =\frac{x^2(1+x^2)}{(1-x)^2(1+x)}.
}{}
It has nonnegative power-series coefficients; its value at $\theta$
is $201617/519840<2/5$. Integration proves the stated finite-tail
bound. If $r_1=4$ then $\ell=2$ or $3$, so the larger mask and the
small path variance belong to the same row; the $r_1\ge5$ row has
the smaller mask.

Take $D=5d/6$, $C=37d/15$ and $\rho_L=1/4$ for every finite $L$.
Substitution in (R30.4.9) gives
\fitdisplay{
\begin{aligned}
 \log\E[\{R_{m,n}P_L^\sigma F\}^8]
 &\le\left[-\frac{128j}{5}-\frac{51}{4}
       +\frac{296}{15}+\frac{21}{4}+\frac{153}{16}\right]d\\
 &=-\left(\frac{128j}{5}-\frac{5231}{240}\right)d
 \le-\frac{913}{240}d<-3d.
 \end{aligned}
}{}
The actual $\chi=\mu_\ell$ was retained before using $0\le\chi\le1$.
This proof also retains $\|V\|_2\le3\sqrt d/4$ on the whole domain.

\paragraph{R30.6 The exact second-shortest length three}
Strict order forces $\ell=2,r_1=3,c=t_1\ge4$. Set
\fitdisplay{
\beta_*=\beta_{c-1},\quad
 R_0=-\beta_*u-B^3G_{c-3}(-B)u,\quad
 \Gamma_0=d+\beta_c+\beta_*.
}{}
The actual pair estimate in Section~\ref{proof:Q5-C181} gives
$f_3+h_c\le2-\Gamma_0+R_0$ and
$\bar f_2f_3h_c\le\bar f_2e^{R_0-\Gamma_0}$, with
$\Gamma_0\ge35d/16$. The $c=4,5$ values are respectively $3d$ and
$3d-\delta_2$; for $c\ge6$ use $d+2\tau\ge35d/16$.
Let
\fitdisplay{
S=\sum_{i=2}^j\sum_{a=r_i}^{t_i-1}K^ag,\quad
 R=R_0+S,\quad \Gamma=\Gamma_0+16(j-1)d/5.
}{}
At $j=1$, $S=0$. Otherwise $r_2\ge c+1\ge5$, and the finite
terminal mask gives
\fitdisplay{
\|S\|_2<3\sqrt d/20,\quad |S|<\sqrt d/4,\quad
 \sum_{i\ge1}(\E B^iS)_+<d/5.
}{}
The actual shortest arm $\bar f_2$ is still present.

For $0<d\le1/8$ the full first row (R30.4.4) has the following
joint bounds from Section~\ref{proof:Q5-C181}:
\fitdisplay{
\begin{array}{c|ccc}
 &D&(D_1)_+&(D_2)_+\\\hline
 j=1&3d/5&d^2&\delta_3\\
 j\ge2&3d/2&19d/120&d/4
 \end{array}
}{}
These bound the complete $\delta$, including its correlated means.
They use coefficient one on $BR$. Also
\fitdisplay{
\|R_0\|_2\le(d+\tfrac65\theta^{3/2})\sqrt d
 <89\sqrt d/200<9\sqrt d/20,
}{}
since $\sqrt{1+\theta}<6/5$ and $\theta^{3/2}<4/15$.
Thus $\|R\|_2<3\sqrt d/5$ for $j\ge2$.
The Taylor conditions follow from
$|R_0|\le d+39\sqrt d/80<1/3$ and
$|R_0+S|\le d+59\sqrt d/80<1/2$;
$\sqrt d<3/8$ gives $197/640<1/3$ and $257/640<1/2$.

For the necessary finite positive-tail sum, write
$G_{c-3}(-z)=U-zV$, so $r_0=-\beta_*+x^2V+zxU$.
At $c=4$, $(U,V)=(1,0)$; otherwise $0\le U,V\le1$.
For every integer $k\ge0$, the four residue classes are
\fitdisplay{
\begin{array}{ll}
 D_{4k}(R_0)=-\int x^{2k+1}U\,d\nu,&
 D_{4k+1}(R_0)=\beta_*\delta_{2k+1}-\int x^{2k+2}V\,d\nu,\\
 D_{4k+2}(R_0)=\int x^{2k+2}U\,d\nu,&
 D_{4k+3}(R_0)=-\beta_*\delta_{2k+2}+\int x^{2k+3}V\,d\nu.
 \end{array}
}{}
Consequently every finite positive subsum is bounded by
\fitdisplay{
\begin{aligned}
 \sum_{i\ge1}(D_i(R_0))_+
 &\le\frac{\beta_*d+\delta_3+\delta_4}{1-\theta^2}\\
 &\le\frac{d^2+\theta^2(1+\theta)d}{1-\theta^2}\\
 &\le\frac{11701}{27360}d<\frac37d.
\end{aligned}
}{}
The last comparison uses $d\le1/8$ and
$7\cdot11701=81907<82080=3\cdot27360$.
For $S$ the corresponding four-class bound is
\fitdisplay{
\frac{x^3+2x^4+x^5}{(1-x^2)^2}
 =\frac{x^3}{(1-x)^2}\le\frac{2197}{11552}<\frac15.
}{}
Thus $\sum_{i\ge1}(D_i(R))_+<22d/35$ by
$(a+b)_+\le a_++b_+$. There is no independence assumption.

\paragraph{R30.7 The complete second-shortest-three payment}
For $j=1$ use $D=3d/5$ and the following finite-length data in
(R30.4.9):
\fitdisplay{
\begin{array}{c|ccc}
 L&C/d&\rho_L&\text{loss coefficient}\\\hline
 1&3/5&9/20&1303/400\\
 2&53/40&3/10&78/25\\
 3&97/60&3/16&17401/3840\\
 \ge4&72/35&3/25&53341/17500
 \end{array}
}{}
The norm bounds use $1$, $\sqrt\theta<2/3$, $\theta<5/12$,
$\theta^{3/2}<4/15$. The scalar bounds use
$(D_1)_+\le d/8$, $(D_2)_+\le\theta^2d<d/6$ and the finite-tail
bound above. The negative base is
$8\Gamma_0+51d/4\ge121d/4$; the four positive costs are respectively
\fitdisplay{
\begin{array}{ll}
 24/5+137/20+6137/400,&53/5+113/20+272/25,\\
 194/15+19/4+2057/256,&576/35+421/100+16337/2500.
 \end{array}
}{}
Subtracting gives the table, all strictly greater than three. The
tight long row has margin $841/17500$ above three. This includes
$c=4,L=1,m=n=3$.

For $j\ge2$ use $D=3d/2$ and
\fitdisplay{
\begin{array}{c|ccc}
 L&C/d&\rho_L&\text{loss at }j=2\\\hline
 1&3/2&3/5&1523/100\\
 2&379/120&2/5&3109/300\\
 3&107/30&1/4&3001/240\\
 \ge4&149/35&4/25&172509/17500
 \end{array}
}{}
Here $8\Gamma+51d/4\ge1117d/20$ at $j=2$, and the positive costs are
\fitdisplay{
\begin{array}{ll}
 12+161/20+2057/100,&379/15+129/20+1377/100,\\
 428/15+21/4+153/16,&1192/35+453/100+18513/2500.
 \end{array}
}{}
All losses exceed nine. Each further pair adds $16d/5$ to $\Gamma$,
while the joint mask, $D$ and finite positive-tail bounds remain the
same. Both low-domain families retain $\|V\|_2\le11\sqrt d/10$.

For $1/8\le d\le1/3$ keep exactly the same direction. The row and
finite $G$ bounds give
\fitdisplay{
|R_0|\le d+39\sqrt d/80\le197d/80,
 \quad \Gamma_0\ge175d/80,
 \quad |S|\le\tfrac34(j-1)d.
}{}
Since $\sqrt d\le3d$ on this domain,
\fitdisplay{
F\le\bar f_2\exp\{[11/40-49(j-1)/20]d\}.
}{}
The positive connector sends $\bar f_2$ to $w_n^\sigma$; the
remaining pair is $Q_{h,H}$. By (R30.3.4) the resulting loss is
$61/20+98(j-1)/5\ge3$. The domains overlap at $d=1/8$.
Together with the structural case above this proves (R30.1.3).

\paragraph{R30.8 Receiver posts and genuine equal pairs}
For $k\ge1$, the actual negative post block has
\fitdisplay{
P_c^-\le e^{R_c-\Gamma_c},\quad
 R_c=\sum_{\nu=1}^k\sum_{a=u_\nu}^{v_\nu-1}B^au=BZ_c,
 \quad\Gamma_c\ge16kd/5.
}{}
Here extracting a factor $B$ is a polynomial identity, not inversion.
The finite mask gives
\fitdisplay{
\|Z_c\|_2^2\le\theta^3 C_*d=2197d/16245<4d/25,
 \quad \E R_c\le\frac{\delta_4}{1-\theta^2}
 \le2197d/27360<d/12.
}{}
The first possible positive mean comes from original power seven.
On the whole domain of \mbox{\hyperref[part:R30.5]{part R30.5}} and the low domain of R30.7,
combine $V+Z_c$ before applying the exponential lemma once. The
extra cost, including all cross terms, is at most
\fitdisplay{
\begin{aligned}
 8\E R_c+34\|V\|_2\|Z_c\|_2+17\|Z_c\|_2^2-8\Gamma_c
 &\le[2/3+374/25+68/25-128k/5]d\\
 &=[1376/75-128k/5]d\le-544d/75<0.
 \end{aligned}
}{}
The terminal and receiver masks may overlap.
On the high domain of R30.7,
$|R_c|\le\|Z_c\|_2<2\sqrt d/5\le6d/5$ gives
$P_c^-\le\exp[-(16k/5-6/5)d]\le1$.
For $k=0$ no additional cost occurs. Finally restore $E_c,E_t\le1$
using positivity of the actual operators. This proves the decorated
bound with every original edge retained.
